% paper.tex — generated by convert_to_tex.py from paper.md.
% DO NOT EDIT THIS FILE BY HAND. Edit paper.md and re-run convert_to_tex.py.
%
% TACL submission, double-blind. Style files in tex/ (tacl2021v1.sty,
% acl_natbib.bst). Compile with: pdflatex paper.tex (twice for refs).

\documentclass[11pt,a4paper]{article}
\usepackage{times,latexsym}
\usepackage{url}
\usepackage[T1]{fontenc}
\usepackage[utf8]{inputenc}
\usepackage{graphicx}
\usepackage{booktabs}
\usepackage{amsmath,amssymb}
\usepackage{textgreek}
\usepackage{newunicodechar}

% Unicode characters that appear in the markdown source. pdflatex doesn't
% have native UTF-8 support for math/Greek glyphs; we map them explicitly.
% If a new Unicode char appears in paper.md and breaks compilation, add a
% \newunicodechar line for it here.
\newunicodechar{ρ}{\ensuremath{\rho}}
\newunicodechar{Δ}{\ensuremath{\Delta}}
\newunicodechar{≈}{\ensuremath{\approx}}
\newunicodechar{≤}{\ensuremath{\leq}}
\newunicodechar{≥}{\ensuremath{\geq}}
\newunicodechar{→}{\ensuremath{\rightarrow}}
\newunicodechar{↔}{\ensuremath{\leftrightarrow}}
\newunicodechar{−}{\ensuremath{-}}
\newunicodechar{×}{\ensuremath{\times}}
\newunicodechar{·}{\ensuremath{\cdot}}
\newunicodechar{²}{\ensuremath{^{2}}}
\newunicodechar{√}{\ensuremath{\surd}}
\newunicodechar{α}{\ensuremath{\alpha}}
\newunicodechar{β}{\ensuremath{\beta}}
\newunicodechar{λ}{\ensuremath{\lambda}}
\newunicodechar{μ}{\ensuremath{\mu}}
\newunicodechar{σ}{\ensuremath{\sigma}}
\newunicodechar{π}{\ensuremath{\pi}}
\newunicodechar{θ}{\ensuremath{\theta}}
\newunicodechar{ε}{\ensuremath{\epsilon}}
\newunicodechar{δ}{\ensuremath{\delta}}
\newunicodechar{γ}{\ensuremath{\gamma}}
\newunicodechar{η}{\ensuremath{\eta}}
\newunicodechar{ω}{\ensuremath{\omega}}
\newunicodechar{Ω}{\ensuremath{\Omega}}
\newunicodechar{Σ}{\ensuremath{\Sigma}}
\newunicodechar{Π}{\ensuremath{\Pi}}
\newunicodechar{Θ}{\ensuremath{\Theta}}
% Section sign (§) is already in T1 fonts.

% Pandoc emits \tightlist for compact lists; define it as a no-op so we
% don't have to enable pandoc's full default preamble.
\providecommand{\tightlist}{\setlength{\itemsep}{0pt}\setlength{\parskip}{0pt}}

% Pandoc emits \pandocbounded around figures with size constraints.
\providecommand{\pandocbounded}[1]{#1}

% TACL style file lives in tex/. Tell LaTeX to look there.
\makeatletter
\def\input@path{{tex/}}
\makeatother
\usepackage[]{tacl2021v1}

% Figures live in figures/.
\graphicspath{{figures/}}

\title{Calibrated Linear Structure of Verbal Probability in Pretrained Sentence Embeddings}

\author{Anonymous Submission \\
  \texttt{anonymized@example.com}}

\date{}

\begin{document}
\maketitle

\begin{abstract}
Epistemic hedging --- words like ``probably,'' ``certainly,'' and ``possibly'' --- has been psychometrically calibrated since the 1960s, with median probability values for verbal expressions remarkably stable across more than fifty years of replication (Mosteller \& Youtz, 1990; Vogel et al., 2022; Wintle et al., 2019). We investigate whether this calibration is recoverable as a \textbf{verbal probability axis} in frozen pretrained pooled sentence embeddings. Within syntactic types, verbal probability is the dominant linear direction in five architecturally diverse pretrained embedding models (Spearman ρ \textgreater{} 0.90 supervised, 0.67--0.95 leave-one-out on the 53-expression Mosteller dataset). The axis cross-validates against two independent psychometric datasets (Vogel 2022 meta-analysis: ρ = 0.991, MAE = 5.1\%; Wintle 2019, n ≈ 924: ρ = 0.967, MAE = 3.6\%), transfers zero-shot to eight typologically diverse languages via a multilingual model (mean ranking correlation \textbar ρ\textbar{} = 0.93), survives 12× dimensional compression on a Matryoshka-trained model (LOO ρ drop = 0.027 at 64 dimensions), and passes permutation null-hypothesis tests (p ≤ 0.005). Rank-1 projection erasure of the axis from cross-type embeddings raises mean absolute prediction error by 18--25 percentage points on the strongest-aligned axis pair (predicative ↔ modal), with degradation correlating with inter-axis cosine alignment at Pearson r = 0.92 (mxbai-embed-large) / 0.86 (qwen3-embedding) across all 12 ordered cross-type pairs. The effect functionally exceeds a cosine-matched random-direction control on both models (MAE z-score \textgreater{} +2 across all four headline pair × model × direction combinations; 9 of 12 ordered pairs functional on both models in the MAE view) --- evidence the axis carries probability content beyond what its geometric overlap with the target axis alone would explain. Recent work has documented analogous structure in two separate model contexts: linear features in \emph{decoder LLM} internal states (Ji et al., 2025; Marks \& Tegmark, 2023), where token-level causal intervention is natural, and verbal-probability calibrations behaviorally elicited from prompted LLMs (Belém et al., 2024; Tang et al., 2026). Pretrained pooled sentence embedding models lack a token-generation pathway, making representation-level intervention the model-class analogue and human psychometric data the appropriate calibration target. The empirical convergence across three psychometric datasets, six embedding models, eight languages, and the functional-erasure result is consistent with distributional embedding geometry and human probability semantics being convergent measurements of one underlying epistemic dimension.
\end{abstract}

\section{Introduction}\label{introduction}

Mikolov et al.~(2013) demonstrated that semantic relationships manifest as consistent vector directions in word embedding space --- \emph{king − man + woman ≈ queen}. The Linear Representation Hypothesis (LRH) generalizes this observation: many human-meaningful concepts correspond to linear directions in language model representations (Park et al., 2024). Three lines of recent work bear on whether epistemic content fits this pattern.

\textbf{Linear features in decoder LLM internal states.} Marks and Tegmark (2023) identify a truth direction in residual stream activations. Bürger et al.~(2024) refine truth representations to handle polarity confounds. Yu et al.~(2025) generalize one-dimensional truth directions to multi-dimensional cones. Ji et al.~(2025) identify a single linear ``verbal uncertainty'' feature in residual streams that can be causally manipulated to reduce overconfident hallucinations. Lepori et al.~(2026) link decoder hidden-state directions to graded human plausibility and modal judgments. Valois et al.~(2025) extend single-token concept directions to multi-token ``frames'' via the Frame Representation Hypothesis, with steering experiments on the unembedding matrix. All of these works share a methodological commitment: they study internal states of decoder LLMs --- residual streams, unembedding matrices, layer-wise activations --- where token-level causal intervention (steering, activation patching) is natural because the model generates tokens.

\textbf{Verbal-probability psychometrics elicited behaviorally from language models.} A parallel literature studies how language models interpret hedge expressions like ``likely,'' ``probably,'' and ``possibly'' by \emph{prompting} the model and comparing its outputs to human numerical judgments. Sileo and Moens (2023) probe pretrained models for words-of-estimative-probability classification (with fine-tuning). Belém et al.~(2024) compare LLM and human numerical interpretations of linguistic uncertainty expressions through prompted-output protocols. Tang et al.~(2026) extend this to multilingual estimative uncertainty. Wang et al.~(2024) calibrate certainty-phrase distributions using transport-based methods. These papers establish that verbal probability semantics are recoverable from language models \emph{behaviorally}, but they do not extract a geometric structure from frozen embedding representations.

\textbf{Pretrained sentence embedding probing for human-aligned graded properties.} A third line studies pooled sentence representations directly. Most directly relevant: Jacobs et al.~(2022) probe mean-pooled RoBERTa sentence embeddings for human cloze-based uncertainty (next-word predictability and sentential constraint), demonstrating that pooled embedding spaces encode some form of human-aligned uncertainty. Chen et al.~(2019) show that fine-tuned BERT-based regressors can predict continuous human probability judgments in sentence-pair NLI. Schuster et al.~(2019) predict human scalar pragmatic ratings from sentence encoders. Pei and Jurgens (2021) regress continuous certainty in scientific writing. These works establish that pooled sentence embeddings carry \emph{some} graded epistemic content --- but for cloze probability, NLI inference strength, or rhetorical certainty rather than verbal probability semantics, and not as a frozen-geometry linear axis calibrated against psychometric data.

\textbf{The gap.} No prior work has demonstrated a \emph{continuously calibrated verbal probability axis} in \emph{frozen pretrained pooled sentence embeddings}, \emph{aligned with human psychometric judgments} and \emph{validated against independent psychometric datasets}. We address this gap. We use the Mosteller and Youtz (1990) dataset, which provides empirically calibrated probability distributions for 53 verbal expressions (e.g., ``likely'' → median 71.1\%, ``possible'' → median 38.5\%) collected from 238 science writers. We embed these in controlled sentence templates, extract difference vectors from bare claims, and test whether the resulting geometry encodes a calibrated probability axis. We further test cross-validation against two independent psychometric datasets (Vogel et al.~2022; Wintle et al.~2019), zero-shot ranking transfer to eight typologically diverse languages, robustness to dimensional truncation, null-hypothesis controls, and \textbf{concept erasure as functional validation} --- the embedding-class analogue of the causal-intervention experiments standard in the decoder-LLM linear-feature literature.

\textbf{Our contributions:}

\begin{enumerate}
\def\labelenumi{\arabic{enumi}.}
\item
  \textbf{A continuously calibrated verbal probability axis in frozen pretrained pooled sentence embeddings.} Within syntactic types, verbal probability is the dominant linear direction across five architecturally diverse pretrained embedding models --- Spearman ρ \textgreater{} 0.90 supervised, 0.67--0.95 leave-one-out on the Mosteller (1990) dataset. This is continuous calibration, not classification: predicted probabilities track Mosteller medians on the same scale humans use. To our knowledge this is the first demonstration of such structure in this model class --- prior work has shown human-aligned uncertainty behavior in prompted LLMs (Belém et al., 2024; Tang et al., 2026), linear uncertainty features in decoder activations (Ji et al., 2025; Lepori et al., 2026), and pooled-embedding probes for cloze-style uncertainty (Jacobs et al., 2022) and NLI inference strength (Chen et al., 2019), but not a psychometrically calibrated verbal probability axis in frozen pooled sentence embeddings.
\item
  \textbf{Three-dataset psychometric triangulation.} The axis cross-validates against two independent psychometric datasets spanning over fifty years of replication. Mosteller (1990) → Vogel (2022) meta-analysis: ρ = 0.991, MAE = 5.1\% (mxbai-embed-large, predicative). Mosteller (1990) → Wintle (2019, n ≈ 924): ρ = 0.967, MAE = 3.6\%. Convergence across datasets collected by independent groups in different decades on different populations argues against an artifact of any single calibration source.
\item
  \textbf{Cross-linguistic ranking transfer.} Using a multilingual embedding model (bge-m3, XLM-RoBERTa, 1024d), an axis trained on English Mosteller data ranks hedge expressions zero-shot in eight typologically diverse languages (German, French, Spanish, Chinese, Japanese, Korean, Arabic, Hindi), mean \textbar ρ\textbar{} = 0.928. Calibration MAE varies by language (Japanese modal MAE = 18\%, Arabic modal MAE = 3.6\%); we treat ranking as the trustworthy transfer claim and calibration as preliminary, pending native-speaker psychometric validation.
\item
  \textbf{Syntactic-type confound.} When expression types are mixed, PCA PC1 captures syntactic frame (28.7\% variance), not probability; the probability signal lives on PC2. Within-type analysis is essential and uses syntax-controlled templates (``It is \{X\} that the experiment will succeed'', ``The experiment will \{X\} succeed'', ``There is a \{X\} that the experiment will succeed''). This finding parallels the polarity confound Bürger et al.~(2024) report for truth directions in decoder LLMs.
\item
  \textbf{Cross-architecture robustness.} The finding holds across BERT (nomic-embed-text v1.5, 768d), BERT-large (mxbai-embed-large, 1024d), Mixture-of-Experts (nomic-v2-moe, 768d), decoder-derived embedding heads (embeddinggemma 300M, 768d; qwen3-embedding, 4096d), and multilingual XLM-R-based (bge-m3, 1024d) --- five architectural families, five training procedures, one phenomenon.
\item
  \textbf{Functional validation via concept erasure.} Rank-1 projection of the probability axis out of sentence embeddings degrades cross-syntactic-type probability prediction \texttt{{[}PLACEHOLDER:\ Δρ,\ MAE{]}} more than a cosine-matched random-direction erasure (z = \texttt{{[}PLACEHOLDER{]}}). The probability axis is functionally carrying probability information used by downstream linear decoders, not merely correlated with embedding geometry.
\item
  \textbf{Null-hypothesis controls.} Permutation tests on Mosteller medians (1000 shuffles) yield p ≤ 0.005 across all four syntactic types. Non-epistemic adjective controls placed in the same templates produce LOO ρ = 0.31 versus 0.87 for real hedges (clean negative on the load-bearing pure-non-epistemic set; we report and address an LOO ρ = 0.69 trend on a mixed set with epistemic leakage). Random Uniform(0,100) labels match real values in ≤ 6 of 1000 trials (p ≤ 0.006).
\item
  \textbf{Informative negative results.} Interpretive precision (Mosteller IQR width) is \emph{not} in the embedding geometry (ρ ≈ 0 across four models) --- embeddings encode a single stable meaning per expression rather than the spread of human judgments. Ensemble combinations across the four within-type axes do not beat the single modal axis on novel-phrase generalization. Negation breaks linear composition (cosine 0.22 vs.~0.82--0.89 for non-negated compounds), which is the geometrically appropriate behavior for a probability-reflecting axis.
\end{enumerate}

The remainder of the paper proceeds as follows. Section 2 positions our work against the linear-feature interpretability literature and the verbal-probability psychometrics literature with structured gap-statements per cited work. Section 3 describes the data, models, axis-extraction procedure, and evaluation protocol. Section 4 presents the eight contributions in order. Section 5 discusses the model-class gap, the bridge between distributional geometry and psychometric semantics, and the four-dimensional epistemic subspace structure. Section 6 names the limitations the paper does not address. Supplementary materials contain full per-model × per-dataset tables, all 8-language per-phrase data, the Wintle verification trail, bootstrap confidence intervals, and concept-erasure z-scores per pair.

\section{Related Work}\label{related-work}

Three lines of recent work motivate, but do not subsume, the question we ask. The first establishes that pretrained language models encode many human-meaningful concepts as linear directions, but does so almost entirely on decoder LLM internal states with token-level causal intervention as the validation method. The second establishes that verbal-probability semantics are recoverable from large language models, but elicits them behaviorally through prompting rather than geometrically from representations. The third establishes that pooled sentence representations carry graded epistemic content, but for cloze probability, NLI inference strength, or rhetorical certainty rather than for the verbal-probability hedge lexicon, and not as a frozen, psychometrically calibrated linear axis. We position our contribution at the intersection these three lines leave open, and close §2 with an explicit gap statement.

\subsection{Linear features in decoder LLM internal states}\label{linear-features-in-decoder-llm-internal-states}

The Linear Representation Hypothesis --- that human-meaningful concepts correspond to linear directions in language model representations --- has been formalized and tested most extensively on decoder LLM internal states (Park et al., 2024). Marks and Tegmark (2023) identify a single ``truth direction'' in residual stream activations that separates true from false factual statements; Bürger et al.~(2024) refine this representation to handle a polarity confound; Yu et al.~(2025) generalize one-dimensional truth directions into multi-dimensional cones for propositional facts. Within this lineage, Ji et al.~(2025) is the most direct framing-level comparator to the present work. They identify a single ``Verbal Uncertainty Feature'' (VUF) in the residual streams of Llama-3.1-8B, Mistral-7B, and Qwen-2.5-7B by computing a difference-in-means direction over question--answer pairs labeled high- versus low-verbal-uncertainty by an LLM-as-a-Judge protocol, and they validate the feature causally by steering activations along it at inference time to produce more or less hedged generations and reduce overconfident hallucinations. Lepori et al.~(2026) extract ``modal difference vectors'' from late-layer hidden states of GPT-2, Llama-3.2, OLMo-2, and Gemma-2 models that linearly discriminate four scenario-level modal categories --- \emph{probable}, \emph{improbable}, \emph{impossible}, \emph{inconceivable} --- across six expert-labeled datasets, demonstrate that projection magnitudes onto these vectors track fine-grained human categorization behavior, and provide preliminary causal-steering evidence that the directions are not epiphenomenal. Valois et al.~(2025) extend the LRH from single tokens to multi-token ``frames'' living on the unembedding matrix and demonstrate functional content via Top-\emph{k} Concept-Guided Decoding. All four of these works share two methodological commitments: they operate on \textbf{decoder LLM internal activations} --- residual streams, late-layer hidden states, unembedding matrices --- and they validate functional content via \textbf{token-level causal intervention} (steering, activation patching, guided decoding), naturally available because the model generates tokens. The lines of evidence they assemble are complementary to ours within the LRH frame; the model class, the available intervention method, and (for Ji and Lepori) the calibration target differ from what we study, in ways §2.5 makes explicit.

\subsection{Verbal-probability semantics elicited behaviorally from language models}\label{verbal-probability-semantics-elicited-behaviorally-from-language-models}

A parallel literature investigates how language models interpret verbal probability expressions like ``likely,'' ``probably,'' and ``highly unlikely'' by \textbf{prompting} the model and comparing its numerical outputs to human survey baselines. Sileo and Moens (2023) probe pretrained models for words-of-estimative-probability classification using fine-tuned, task-specific objectives. Belém et al.~(2024) compare ten LLMs against 94 human participants on 14 uncertainty expressions drawn from Wallsten et al.~(1986a, 2008), framing the task as theory-of-mind (``Sonia believes it is unlikely it will rain today'' --- quantify Sonia's certainty, not your own) and finding that 7 of 10 LLMs map uncertainty to probabilities in human-aligned ways but are substantially more biased by prior knowledge than humans are. Tang et al.~(2026) evaluate 17 Words of Estimative Probability across GPT-3.5, GPT-4, two Llama-2 variants, and ERNIE-4 against a human survey baseline, finding human--LLM alignment for high-certainty expressions but divergence on the majority of mid-range expressions and additional divergence under gendered-role and Chinese-language prompting. Wang et al.~(2024) calibrate certainty-phrase distributions against human data using transport-based methods at the phrase level, again without embedding analysis. These works collectively establish that verbal probability semantics are quantitatively recoverable from contemporary LLMs and that human population-level probability mappings for hedge expressions are robust enough to serve as a comparison standard. They locate the phenomenon \textbf{behaviorally}, in model outputs, rather than geometrically, in model representations.

\subsection{Probing pretrained sentence representations for human-aligned graded properties}\label{probing-pretrained-sentence-representations-for-human-aligned-graded-properties}

A third, smaller line studies pooled or encoded sentence representations directly. The strongest precursor in our model class is Jacobs et al.~(2022): mean-pooled RoBERTa sentence embeddings, with the critical final word replaced by \texttt{\textless{}mask\textgreater{}}, are passed through 282 leave-one-out regularized logistic-regression classifiers trained to discriminate strong- versus weak-sentential-constraint sentences from the Federmeier et al.~(2007) stimulus set; the classifier's predicted probability is then correlated against the cloze probability of the masked critical word (\(\hat\rho = 0.43\)) and against the cloze entropy of human responses (\(\hat\rho = -0.32\)). The differentiation runs two layers deep. \emph{Phenomenon:} their target is sentential constraint operationalized through next-word predictability --- uncertainty about which word will come next given a preceding context --- whereas our target is the verbal-probability semantics of explicit hedge expressions (``probably'' → 71\%, ``possible'' → 39\%) collected from human raters. \emph{Method:} they train per-item binary classifiers and then correlate classifier output with a continuous human signal external to the classifier; we directly regress pooled sentence embedding differences onto continuous psychometric medians, extracting a single linear axis whose calibration is the geometric object of study. They do not extract a calibrated verbal-probability axis, do not cross-validate against independent psychometric datasets, and do not test cross-linguistic transfer. Chen et al.~(2019) refine NLI annotation from categorical labels to continuous subjective probability via the u-SNLI dataset (61,597 MTurk-elicited probability annotations on premise--hypothesis pairs), then fine-tune scalar regression models from BERT-derived sentence representations; the phenomenon is sentence-pair entailment probability, the method is fine-tuned scalar regression rather than frozen-geometry probing, and no general probability axis is identified. Schuster et al.~(2019) predict human scalar pragmatic ratings (scalar implicature strength) from task-trained sentence encoders. Pei and Jurgens (2021) regress continuous certainty in scientific writing from fine-tuned transformer encoders, locating rhetorical certainty rather than verbal-probability semantics. Lee et al.~(2015) and Stanovsky et al.~(2017) predict scalar factuality with structured task-specific models. Bhatia (2016) shows distributional vector-space models predict subjective probability judgments for real-world events from keyword associations, predating sentence-embedding architectures. Shen et al.~(2023) layer model and data uncertainty over pretrained sentence encoders, modeling distributional embedding noise rather than human psychometric semantics. Together, these works establish that pooled sentence representations carry \emph{some} graded epistemic content --- but for cloze probability, NLI inference strength, scalar implicature, rhetorical certainty, factuality, or generic event likelihood, and not as a frozen, psychometrically calibrated verbal-probability axis cross-validated against independent psychometric datasets.

\subsection{Verbal-probability psychometrics as ground truth}\label{verbal-probability-psychometrics-as-ground-truth}

The data side of this paper draws on a sixty-year tradition of verbal-probability calibration studies. Lichtenstein and Newman (1967) and Beyth-Marom (1982) established that population-level mappings from verbal expressions to numerical probabilities are systematic and reproducible. Wallsten et al.~(1986a, 2008) and Budescu and Wallsten (1995) developed the methodological framework --- non-verifiable contexts, careful statement construction, distributional rather than point estimates --- that the modern literature continues to use. Mosteller and Youtz (1990) provide the primary calibration corpus we train on: 53 verbal probability expressions rated by 238 science writers, with median, P25, P75, and IQR per expression. Vogel et al.~(2022) is a systematic review and meta-analysis of 21 such studies spanning 1967--2018, which we use as an independent cross-validation target. Wintle et al.~(2019, \(n \approx 924\)) is a third independent dataset, collected on a different population with a different protocol, providing an additional cross-validation target. Lassiter (2017) develops the formal-semantics theory of graded modality --- degree-based rather than possible-worlds --- that motivates treating ``probably,'' ``possibly,'' and their kin as continuous-valued operators in the first place. Schockaert (2022) provides a formal framework for embeddings as epistemic states, supplying the measurement-theoretic scaffolding for the convergent-measurement perspective we adopt in §5.

\subsection{The gap}\label{the-gap}

Across these three lines, no prior work demonstrates the conjunction this paper fills:

\begin{enumerate}
\def\labelenumi{\arabic{enumi}.}
\tightlist
\item
  \textbf{Continuously calibrated} rather than categorically classified. Sileo and Moens (2023) classify; Jacobs et al.~(2022) train binary classifiers and correlate their outputs with continuous external signals; we regress pooled embedding differences directly onto psychometric medians, on the same percentage scale humans use.
\item
  \textbf{Verbal probability} specifically --- the hedge lexicon (``likely,'' ``probably,'' ``doubtful,'' ``almost certain'') --- rather than truth, factuality, sentential constraint, NLI inference probability, modal categorization of event plausibility, or scalar implicature strength.
\item
  \textbf{Frozen pretrained pooled sentence embeddings}, rather than decoder LLM internal activations (Ji et al., 2025; Lepori et al., 2026; Marks and Tegmark, 2023; Yu et al., 2025; Valois et al., 2025) or fine-tuned task-specific models (Sileo and Moens, 2023; Chen et al., 2019; Pei and Jurgens, 2021).
\item
  \textbf{Aligned with human psychometric judgments} (Mosteller and Youtz, 1990; Wallsten et al., 1986a) rather than with model-internal LLM-as-a-Judge labels (Ji et al., 2025) or expert-assigned modal-category labels (Lepori et al., 2026).
\item
  \textbf{Validated against independent psychometric datasets} (Vogel et al., 2022; Wintle et al., 2019), spanning over fifty years of replication, rather than against a single calibration source.
\end{enumerate}

Each qualifier on its own excludes a class of near-misses; together they identify the open territory this paper occupies. The contribution is parallel in spirit to the decoder-LLM linear-feature literature --- same family of finding, namely that a graded epistemic property is recoverable as a linear direction in a pretrained representation --- but is conducted in a model class where token-level causal intervention is not available, where representation-level concept erasure (§4.4) is the appropriate analogue, and where human psychometric data, rather than model behavior, is the appropriate calibration target.

\section{Method}\label{method}

\subsection{Data}\label{data}

\subsection{Models}\label{models}

\subsection{Axis extraction}\label{axis-extraction}

\subsection{Evaluation}\label{evaluation}

\section{Results}\label{results}

\subsection{The syntactic confound and within-type probability axes}\label{the-syntactic-confound-and-within-type-probability-axes}

\subsection{Three-dataset psychometric triangulation}\label{three-dataset-psychometric-triangulation}

\subsection{Cross-linguistic ranking transfer}\label{cross-linguistic-ranking-transfer}

\subsection{Functional validation via concept erasure}\label{functional-validation-via-concept-erasure}

\subsection{Robustness: dimensional truncation, novel-phrase generalization, compound composition, bimodality probe}\label{robustness-dimensional-truncation-novel-phrase-generalization-compound-composition-bimodality-probe}

\subsection{Null-hypothesis controls}\label{null-hypothesis-controls}

\subsection{Negative results}\label{negative-results}

\section{Discussion}\label{discussion}

\section{Limitations}\label{limitations}

\section{Conclusion}\label{conclusion}

% References: TACL uses acl_natbib. A refs.bib file is not yet built; the
% draft still has citations in (Author et al., Year) markdown form. When
% refs.bib lands, uncomment the two lines below and re-run.
%
% \bibliography{refs}
% \bibliographystyle{tex/acl_natbib}

\end{document}
